
\documentclass[a4paper,12pt]{article}
\usepackage[utf8]{inputenc}
\usepackage[russian]{babel}
\usepackage{amsmath, amsfonts, amssymb, amsthm}
\usepackage{graphicx}
\usepackage{hyperref}
\usepackage{geometry}
\usepackage{float}
\geometry{left=2cm,right=2cm,top=2cm,bottom=2cm}

\title{\textbf{Лабораторная работа №3}:\\
Трек 7. Сглаживание Нестерова и Проксимальные методы}
\author{Выполнили: Панов Богдан, Кахниашвили Иоанн, Кретов Дмитрий}

\begin{document}

\maketitle
\newpage

\section*{Постановка задачи}

В данной работе исследуются методы решения задачи минимизации с L1-регуляризацией:
\[
\min_{x \in \mathbb{R}^n} \; \phi(x) := f(x) + \lambda \|x\|_1,
\]
где $f(x)$ — гладкая выпуклая функция потерь (пуассоновская регрессия или логистическая классификация из пакета №2), $\lambda > 0$ — коэффициент регуляризации.

\subsection*{Функции потерь (пакет №2)}

\textbf{Пуассоновская регрессия:}
\[
f(x) = \frac{1}{m}\mathbf{1}^\top(\exp(Ax) - y \odot Ax) + \frac{\mu}{2}\|x\|_2^2, \qquad
\nabla f(x) = \frac{1}{m}A^\top(\exp(Ax) - y) + \mu x
\]

\textbf{Логистическая классификация:}
\[
f(x) = \frac{1}{m}\mathbf{1}^\top\log(1 + \exp(-y \odot Ax)) + \frac{\mu}{2}\|x\|_2^2, \qquad
\nabla f(x) = \frac{1}{m}A^\top(-y \odot \sigma(-y \odot Ax)) + \mu x
\]

Рассматриваемые методы:
\begin{itemize}
    \item \textbf{Субградиентный метод} — применим напрямую к $\phi(x)$, не требует дифференцируемости
    \item \textbf{ISTA} (проксимальный градиентный метод) — использует структуру $\phi = f + h$
    \item \textbf{FISTA} (ускоренный метод Нестерова) — добавляет импульс к ISTA
    \item \textbf{Метод Франка-Вулфа} — минимизация $f(x)$ на L1-шаре $\|x\|_1 \le R$
    \item \textbf{Барьерный метод} — решает эквивалентную задачу с ограничениями
\end{itemize}

\section*{Описание методов}

\subsection*{Субградиентный метод}

На каждой итерации выбирается произвольный субградиент $g_k \in \partial \phi(x_k)$ и делается шаг:
\[
x_{k+1} = x_k - \alpha_k g_k, \qquad \alpha_k = \frac{\alpha_0}{\sqrt{k+1}}.
\]
Субградиент вычисляется как $g_k = \nabla f(x_k) + \lambda \operatorname{sign}(x_k)$.
Теоретическая скорость сходимости составляет $O(1/\sqrt{k})$, что является оптимальным для несглаженных выпуклых задач.

\subsection*{ISTA (проксимальный градиентный метод)}

Представляет задачу в виде $\phi = f + h$, где $h(x) = \lambda\|x\|_1$ — простая часть с известным
проксимальным оператором (мягкое пороговое преобразование):
\[
\operatorname{prox}_{\alpha h}(x)_i = \operatorname{sign}(x_i)\max(|x_i| - \alpha\lambda,\; 0).
\]
Шаг метода:
\[
x_{k+1} = \operatorname{prox}_{h/L_k}\!\bigl(x_k - \tfrac{1}{L_k}\nabla f(x_k)\bigr),
\]
где $L_k$ — оценка константы Липшица градиента, подбираемая бэктрекингом. Скорость сходимости $O(1/k)$.

\subsection*{FISTA (ускоренный метод Нестерова)}

Добавляет к ISTA инерционный шаг:
\[
y_{k+1} = \operatorname{prox}_{h/L_k}\!\bigl(y_k - \tfrac{1}{L_k}\nabla f(y_k)\bigr), \quad
t_{k+1} = \frac{1 + \sqrt{1 + 4t_k^2}}{2}, \quad
x_{k+1} = y_{k+1} + \frac{t_k - 1}{t_{k+1}}(y_{k+1} - y_k).
\]
Скорость сходимости $O(1/k^2)$ — квадратичное ускорение без дополнительных вычислений оракула.

\subsection*{Метод Франка-Вулфа (условный градиент)}

Решает задачу $\min_{x:\|x\|_1 \le R} f(x)$. На каждом шаге решает задачу линейной минимизации:
\[
y_k = \operatorname{argmin}_{\|y\|_1 \le R} \langle \nabla f(x_k), y \rangle = -R \cdot \operatorname{sign}\bigl(\nabla f(x_k)[i^*]\bigr) \cdot e_{i^*},
\quad i^* = \operatorname{argmax}_i |\nabla f(x_k)_i|,
\]
затем обновляет $x_{k+1} = x_k + \gamma_k (y_k - x_k)$.
Критерий остановки — зазор Франка-Вулфа: $\langle \nabla f(x_k), x_k - y_k \rangle \le \varepsilon$.
Рассматривались два варианта выбора шага: $\gamma_k = 2/(k+2)$ (стандартный) и Армихо (линейный поиск).

\subsection*{Барьерный метод}

Переформулирует задачу с L1-регуляризацией через эквивалентную задачу с ограничениями:
\[
\min_{x, u} \; f(x) + \lambda \mathbf{1}^\top u \quad \text{s.t.} \quad -u_i \le x_i \le u_i, \; \forall i,
\]
и применяет логарифмический барьер:
\[
F_t(x, u) = t\bigl(f(x) + \lambda \mathbf{1}^\top u\bigr) - \sum_i \ln(u_i - x_i) - \sum_i \ln(u_i + x_i).
\]
Внешний цикл увеличивает $t \leftarrow \mu t$ ($\mu = 5$), внутренний — метод Ньютона по $(x, u)$.
Критерий остановки: $2n/t \le \varepsilon$.

\newpage
\section*{Эксперимент 1: Субградиент vs ISTA vs FISTA}

\subsection*{Постановка}

Сравниваются три метода на задачах $\min_{x} f(x) + \lambda\|x\|_1$
с $\lambda = 10^{-3}$ на датасетах \texttt{space\_ga} (пуассоновская регрессия,
$m = 3107$, $n = 6$) и \texttt{phishing} (логистическая классификация,
$m = 11055$, $n = 68$). Начальная точка $x_0 = 0$.
Метрика — относительная невязка:
\[
\frac{\phi(x_k) - \phi^*}{\phi(x_0) - \phi^*}.
\]

\subsection*{Результаты}

\begin{figure}[H]
    \centering
    \includegraphics[width=0.95\textwidth]{img/regression_convergence.png}
    \caption{Сходимость субградиентного метода, ISTA и FISTA на задаче пуассоновской регрессии (\texttt{space\_ga}, $\lambda=10^{-3}$). Слева --- невязка от числа итераций, справа --- от времени.}
    \label{fig:reg_conv}
\end{figure}

\begin{figure}[H]
    \centering
    \includegraphics[width=0.95\textwidth]{img/classification_convergence.png}
    \caption{Сходимость субградиентного метода, ISTA и FISTA на задаче логистической классификации (\texttt{phishing}, $\lambda=10^{-3}$). Слева --- невязка от числа итераций, справа --- от времени.}
    \label{fig:cls_conv}
\end{figure}

\subsection*{Наблюдения}

\begin{itemize}
    \item \textbf{Субградиентный метод} практически не прогрессирует: невязка остаётся на уровне $10^{-1}$ даже после 3000 итераций. Это соответствует теоретической скорости $O(1/\sqrt{k})$: при $k=3000$ выигрыш составляет лишь $\sim 18$ раз. При $\lambda=10^{-3}$ задача хорошо структурирована, и отсутствие эксплуатации L1-структуры становится серьёзным недостатком.

    \item \textbf{ISTA} монотонно убывает от $10^{-1}$ до $10^{-3}$--$10^{-4}$ за 3000 итераций. Скорость $O(1/k)$ наглядно видна на графике.

    \item \textbf{FISTA} демонстрирует характерные нарастающие осцилляции с последующим резким спадом до уровня $10^{-14}$--$10^{-16}$ (машинный ноль) уже через $600$--$1000$ итераций в задаче регрессии и $\sim 1100$ итераций в задаче классификации. Это ускорение в $\sim 3$--$5$ раз по числу итераций и примерно в такое же число раз по времени.
\end{itemize}

\textbf{Вывод:} субградиентный метод применим, но существенно уступает проксимальным методам, которые эксплуатируют L1-структуру через оператор мягкого порогования. FISTA обеспечивает квадратичное ускорение по сравнению с ISTA без дополнительных обращений к оракулу.

\newpage
\section*{Эксперимент 2: Теоретические скорости ISTA и FISTA}

\subsection*{Постановка}

На логарифмическом графике (log-log) сравниваются эмпирические кривые сходимости ISTA и FISTA
с теоретическими Reference-кривыми $O(1/k)$ и $O(1/k^2)$.

\begin{figure}[H]
    \centering
    \includegraphics[width=0.95\textwidth]{img/ista_fista_rates.png}
    \caption{Скорости сходимости ISTA и FISTA на log-log графике с теоретическими кривыми $O(1/k)$ и $O(1/k^2)$. Слева --- регрессия (\texttt{space\_ga}), справа --- классификация (\texttt{phishing}).}
    \label{fig:rates}
\end{figure}

\subsection*{Наблюдения}

\begin{itemize}
    \item \textbf{ISTA} следует эмпирическому наклону, близкому к $O(1/k)$: на обоих датасетах прямолинейный спад на log-log графике совпадает с теоретической пунктирной кривой.

    \item \textbf{FISTA} на начальных итерациях ($k < 100$) убывает медленнее из-за накопления импульса, затем резко ускоряется и обгоняет ISTA. Осциллирующая форма кривой обусловлена инерционным членом $(t_k - 1)/t_{k+1}$, который на каждом шаге может временно увеличивать значение $\phi(x_k)$. Финальный наклон FISTA круче $O(1/k^2)$, что соответствует теоретической гарантии.

    \item На задаче классификации (\texttt{phishing}) оба метода сходятся медленнее за счёт большей размерности ($n = 68$), но качественная картина сохраняется.
\end{itemize}

\textbf{Теоретический результат (Beck \& Teboulle, 2009):}
Для FISTA выполняется
\[
\phi(x_k) - \phi^* \le \frac{2L\|x_0 - x^*\|_2^2}{(k+1)^2},
\]
что при фиксированных $L$ и $\|x_0 - x^*\|$ даёт улучшение по отношению к ISTA в $O(k)$ раз.

\newpage
\section*{Эксперимент 3: Метод Франка-Вулфа — стандартный шаг vs Армихо}

\subsection*{Постановка}

Минимизируется функция пуассоновской регрессии на L1-шаре $\|x\|_1 \le R = 1.0$
(без L2-регуляризации, чтобы ограничение было активным).
Сравниваются два варианта выбора шага:
\begin{itemize}
    \item Стандартный: $\gamma_k = 2/(k+2)$
    \item Армихо: бэктрекинг с условием $f(x_k + \gamma d_k) \le f(x_k) + c_1 \gamma \nabla f(x_k)^\top d_k$
\end{itemize}

\begin{figure}[H]
    \centering
    \includegraphics[width=0.95\textwidth]{img/frank_wolfe_regression.png}
    \caption{Зазор Франка-Вулфа $\langle \nabla f(x_k),\, x_k - y_k \rangle$ для двух стратегий выбора шага на задаче регрессии (\texttt{space\_ga}, $R=1.0$). Слева --- от итераций, справа --- от времени.}
    \label{fig:fw}
\end{figure}

\subsection*{Наблюдения}

\begin{itemize}
    \item Оба варианта демонстрируют убывание FW-зазора, соответствующее теоретической скорости $O(1/k)$.

    \item \textbf{Armijo} сходится быстрее по числу итераций: уже к 200–500-й итерации зазор ниже, чем у стандартного шага. Это связано с тем, что на начальных итерациях стандартный шаг $2/(k+2)$ мал и недооценивает оптимальное значение.

    \item По \textbf{времени} обе стратегии почти эквивалентны: Armijo тратит дополнительное время на поиск шага, что компенсирует выигрыш в итерациях. По времени до $\varepsilon = 10^{-3}$ оба метода достигают цели примерно одновременно.

    \item На обоих графиках видны осцилляции FW-зазора: поскольку метод Франка-Вулфа ищет экстремальный угол конуса (один ненулевой компонент), функция неглад\-ко убывает вдоль траектории.

    \item Метод сходится к точке на границе L1-шара ($\|x^*\|_1 = R = 1.0$), что подтверждает корректность работы LMO.
\end{itemize}

\textbf{Вывод:} при одинаковой теоретической гарантии $O(1/k)$ вариант с Armijo даёт более равномерную монотонную сходимость, тогда как стандартный шаг проще в реализации и сравним по времени.

\newpage
\section*{Эксперимент 4: Барьерный метод vs ISTA и FISTA}

\subsection*{Постановка}

Для задачи малой размерности ($n = 6$, подвыборка \texttt{space\_ga}: $m = 200$ объектов)
сравниваются барьерный метод, ISTA и FISTA на задаче $\min_x f(x) + 0.05\|x\|_1$.

Барьерный метод работает с 2n-мерным вектором $z = [x, u]$ и вычисляет полный гессиан
$\nabla^2 F_t(z) \in \mathbb{R}^{2n \times 2n}$.

\begin{figure}[H]
    \centering
    \includegraphics[width=0.95\textwidth]{img/barrier_vs_proximal.png}
    \caption{Сравнение ISTA, FISTA и барьерного метода ($n=6$, $\lambda=0.05$, $m=200$). Слева --- от итераций, справа --- от времени.}
    \label{fig:barrier}
\end{figure}

\subsection*{Наблюдения}

\begin{itemize}
    \item \textbf{ISTA и FISTA} сходятся до уровня машинной точности ($\sim 10^{-17}$) за $\sim 50$--$100$ итераций. Для задачи с $n = 6$ разница между ними незначительна по числу итераций, но обе кривые практически совпадают и уходят в ноль стремительно.

    \item \textbf{Барьерный метод} демонстрирует медленную сходимость: относительная невязка остаётся на уровне $\sim 10^{-1}$ на протяжении всего запуска. Это объясняется тем, что при $\mu = 5$ и $t_0 = 1$ барьерный метод находит решение с точностью $2n/t_k$ на каждом внешнем шаге, но для достижения высокой точности требуется $t \gg 2n/\varepsilon = 120$. С другой стороны, большое $\mu$ означает быстрый рост $t$, и через $\sim 8$ итераций $t = 5^8 \approx 4 \cdot 10^5$, что должно давать $2n/t \approx 3 \cdot 10^{-5}$. Наблюдаемая константная невязка указывает на то, что \textbf{барьерный метод сходится к иной точке}, чем ISTA, — это связано с тем, что барьерная функция приближает L1-регуляризацию лишь при $t \to \infty$, и при конечном $t$ находится приближённое решение.

    \item По \textbf{времени} барьерный метод тратит $\sim 2.5$ с (из-за решения системы $2n \times 2n$ на каждой внутренней итерации), тогда как ISTA/FISTA завершают работу за $\sim 0.001$ с.
\end{itemize}

\textbf{Вывод:} барьерный метод применим только для задач малой размерности и не конкурентоспособен с проксимальными методами по скорости. Его преимущество — теоретическая точность решения через эквивалентную LP-форму. Для $n \sim 10^3$ вычислительная стоимость $O(n^3)$ на итерацию делает метод непрактичным.

\newpage
\section*{Эксперимент 5: Влияние $\lambda$ на разреженность решения}

\subsection*{Постановка}

FISTA запускается для набора значений $\lambda \in \{10^{-4}, 5 \cdot 10^{-4}, 10^{-3}, \ldots, 0.5\}$
на полных датасетах \texttt{space\_ga} и \texttt{phishing}.
Измеряется доля нулевых компонент $|\{i : |x^*_i| < 10^{-6}\}| / n$.

\begin{figure}[H]
    \centering
    \includegraphics[width=0.95\textwidth]{img/sparsity_vs_lambda.png}
    \caption{Слева --- разреженность решения (доля нулевых компонент) от $\lambda$ для FISTA. Справа --- значение целевой функции от $\lambda$. Логарифмическая шкала по горизонтали.}
    \label{fig:sparsity}
\end{figure}

\subsection*{Наблюдения}

\begin{itemize}
    \item \textbf{Разреженность монотонно растёт} с $\lambda$: при $\lambda = 10^{-4}$ решение плотное,
    при $\lambda = 0.5$ — полностью нулевое (100\%). Это прямое следствие оператора мягкого порогования:
    компонента $x_i^*$ обнуляется, когда $|\nabla f(x^*)_i| \le \lambda$.

    \item \textbf{Классификация} (\texttt{phishing}) даёт более разреженные решения при одинаковом $\lambda$:
    уже при $\lambda = 5 \cdot 10^{-4}$ доля нулей $\approx 22\%$ (против 0\% для регрессии). Это связано
    с меньшим масштабом градиента логистических потерь (ограничен $\frac{1}{4m}\|A^\top A\|_2$).

    \item \textbf{Значение функции} для регрессии практически не меняется в диапазоне $\lambda \in [10^{-4}, 0.5]$
    — кривая плоская. Это объясняется тем, что при $n = 6$ признаки пространственно ортогональны
    (после нормализации), и регуляризация лишь незначительно сдвигает оптимум.
    Для классификации при $n = 68$ рост $\lambda$ заметно увеличивает значение $\phi(x^*)$.
\end{itemize}

\textbf{Вывод:} L1-регуляризация является мощным инструментом отбора признаков. Порог $\lambda^*$, при котором начинается разреженность, зависит от масштаба градиента и пространственной структуры данных.

\newpage
\section*{Эксперимент 6: Чувствительность субградиентного метода к $\alpha_0$}

\subsection*{Постановка}

Субградиентный метод запускается на задаче логистической классификации (\texttt{phishing})
с различными значениями начального шага $\alpha_0 \in \{0.01, 0.1, 0.5, 1.0, 2.0\}$,
шаг $\alpha_k = \alpha_0 / \sqrt{k+1}$.

\begin{figure}[H]
    \centering
    \includegraphics[width=0.85\textwidth]{img/subgrad_alpha_sensitivity.png}
    \caption{Влияние начального шага $\alpha_0$ на сходимость субградиентного метода (задача классификации \texttt{phishing}).}
    \label{fig:alpha}
\end{figure}

\subsection*{Наблюдения}

\begin{itemize}
    \item При $\alpha_0 \in \{0.01, 0.1\}$ метод почти не прогрессирует за 2000 итераций:
    невязка остаётся на уровне $10^{-1}$--$10^{-4}$. Шаг слишком мал для движения к решению.

    \item При $\alpha_0 \in \{0.5, 1.0, 2.0\}$ метод сходится гораздо быстрее: уже через
    $\sim 50$--$200$ итераций невязка падает до уровня $10^{-14}$. Это объясняется тем, что
    при достаточно большом начальном шаге $\alpha_k$ метод быстро ``находит'' разреженную структуру
    решения (у \texttt{phishing} при $\lambda = 10^{-3}$ часть компонент уходит в ноль),
    после чего субградиент в нуле равен нулю и метод эффективно останавливается.

    \item Более крупный $\alpha_0 = 2.0$ оказывается ненамного лучше $\alpha_0 = 1.0$: кривые
    перекрываются.

    \item Выбор $\alpha_0$ критически важен для субградиентного метода: оптимальный шаг
    зависит от неизвестной константы Липшица и расстояния до оптимума, поэтому в отличие
    от ISTA/FISTA (где $L_k$ подбирается автоматически) субградиентный метод требует ручного тюнинга.
\end{itemize}

\textbf{Вывод:} субградиентный метод с убывающим шагом чрезвычайно чувствителен к выбору $\alpha_0$.
При удачном шаге метод быстро обнаруживает разреженное решение, при неудачном — практически не сходится
в разумное число итераций.

\newpage
\section*{Итоговые выводы}

\begin{enumerate}
    \item \textbf{Проксимальные методы (ISTA, FISTA) принципиально превосходят субградиентный метод}
    при наличии L1-структуры. Эксплуатация оператора мягкого порогования даёт выигрыш в десятки и
    сотни раз по числу итераций и по времени.

    \item \textbf{Ускорение Нестерова (FISTA vs ISTA)} даёт практическое ускорение в $3$--$10$ раз
    по итерациям. Теоретическая гарантия $O(1/k^2)$ подтверждается на обоих датасетах:
    на log-log графике наклон FISTA заметно круче, чем у ISTA.

    \item \textbf{Метод Франка-Вулфа} подходит для L1-ограниченной задачи (не L1-регуляризованной).
    Шаг Армихо ускоряет сходимость по итерациям, но не даёт выигрыша по времени.
    Скорость $O(1/k)$ совпадает с ISTA.

    \item \textbf{Барьерный метод} обеспечивает теоретически точное решение через эквивалентную
    LP-постановку, но требует $O(n^3)$ на итерацию и не конкурентоспособен с проксимальными
    методами даже при $n = 6$. Применим только для малых задач, где важна высокая точность
    и допустима высокая вычислительная стоимость.

    \item \textbf{L1-регуляризация как отбор признаков:} увеличение $\lambda$ монотонно увеличивает
    разреженность решения. Классификационные задачи дают более разреженные решения при том же $\lambda$
    из-за меньшего масштаба градиента логистических потерь.

    \item \textbf{Субградиентный метод} чувствителен к $\alpha_0$ и при удачном выборе может
    быстро найти разреженное решение, однако не имеет адаптивного механизма подбора шага
    в отличие от ISTA/FISTA.
\end{enumerate}

\begin{table}[H]
\centering
\begin{tabular}{|l|c|c|c|c|}
\hline
\textbf{Метод} & \textbf{Скорость} & \textbf{Oracle} & \textbf{Память} & \textbf{Применимость} \\
\hline
Субградиент & $O(1/\sqrt{k})$ & $\nabla \phi$ & $O(n)$ & несглаженные задачи \\
ISTA & $O(1/k)$ & $\nabla f$ + prox & $O(n)$ & $\phi = f + h$ \\
FISTA & $O(1/k^2)$ & $\nabla f$ + prox & $O(n)$ & $\phi = f + h$ \\
Франк-Вулф & $O(1/k)$ & $\nabla f$ + LMO & $O(n)$ & L1-ограничение \\
Барьерный & $O(\log(n/\varepsilon))$ & $\nabla f$, $\nabla^2 f$ & $O(n^2)$ & малое $n$ \\
\hline
\end{tabular}
\caption{Сравнение характеристик исследованных методов.}
\end{table}

\end{document}
